\subsection{Conventions and inherited hypotheses}
\label{sec:v4-inherited}

Retain the field and curve of Section~\ref{sec:degree8-original}:
\[
 K=\Q(\sqrt2),\qquad d=3-\sqrt2,\qquad D=d^2,
 \qquad E:\ y^2=x^3-Dx.
\]
For an affine point with nonzero abscissa write $f(P)=1/x(P)$ and
$p(P)=y(P)/x(P)$. The rational function $z=-x/y$ extends with value zero
at $O$ and $T_0=(0,0)$; its other poles are $(d,0)$ and $(-d,0)$.
The letter $J$ denotes a common comparison point. We write $\jmath$ for
$1/z(J)$, to distinguish that field element from a point or an integer index.
The implementation conjugates the complete alternate-comparison identity:
its boundary value is $1/\overline{z(J)}$, and its primary and target
coefficients in that identity are conjugated as well.

The following arguments retain all inherited qualifications of
Section~\ref{sec:degree8-original} and Appendix~\ref{app:sources}. In particular, the restricted converse, rank and original
arithmetic assertions retain their stated proof status. The prime-pattern
input retains its corrected author-version and erratum conventions,
including the stated version of arXiv:1009.3998v5. The Pan interface retains
the inspected author theorem and the cited Calegari formulation, without
claiming comparison with inaccessible journal text. The Yuan input retains
its author-attested translation qualification. The stable Faltings metric,
arithmetic-Frobenius, diamond-quotient and final-twist conventions are
unchanged. None of the finite source or polynomial checks strengthens these
assumptions.

In particular we use, in exactly their previous scopes, the generator and
torsion statements $2E(K)=\Z P_*$ and $E(K)_{\mathrm{tors}}=E[2](K)$;
the formal-group logarithm and valuation statements; the fixed-system
prime-pattern theorem with arbitrary prescribed finite exclusions; the
height, scalar, anchor and Lucas lemmas; and the restricted-converse and
simultaneous-realization inputs used for the pole predicate. The one
additional explicit trace application below is identified in
Lemma~\ref{lem:quartic-anchor-four}.

Put
\[
 B_* =\prod_{p\le76176}p,\qquad
 H=2(2^{109}210^{1800})B_*^{19}+128,
 \qquad \log_2 H<2093966.
\]
This is $H_{19}$ of Theorem~\ref{thm:quartic-height19}. The
source parameters remain high order $40$, beta order $12$, scalar orders
$(1,1,1,1,32,32)$ and anchor constant $C=3^{2^{25}}$.

We also retain Lemma~\ref{lem:quartic-reciprocal-log}. Explicitly,
write
\[
 \sqrt{1-w}\sum_{j\ge0}\frac{\binom{2j}{j}}{4^j(4j+1)}w^j
       =\sum_{j\ge0}a_jw^j,\qquad
 G_m(w)=\sum_{j=0}^m a_jw^j.
\]
The polynomials used here are $G_9$ and $G_2$. They are nonzero at every
argument $Df^2$ with $f\in K$, by the previously proved dyadic argument.
For a nonidentity formal-group point at either place above a split prime
$p>76176$, the approximation $-G_9(Df^2)/p(P)$ has relative order at least
$40$, and $-G_2(Df^2)/p(P)$ has relative order at least $12$. Actual formal
membership is required: small $z$ alone is not the premise of this lemma.

\subsubsection{The retained scalar, anchor and beta interfaces}
The scalar input is still $\beta=X^2+2$. The scalar batch tests the two
rational coefficients each of $f(J_0),f(S),z(S)/z(J_0)$, with orders
$(1,1,1,1,32,32)$ in that order. Its original pole predicate is retained.
The anchor constant is computed by
\[
 c_0=3,\qquad c_{i+1}=c_i^2\quad(0\le i<25),\qquad C=c_{25}.
\]
Thus the integer $C=3^{2^{25}}$ is a witness value, not a polynomial
coefficient. The four-witness anchor below consumes this same terminal
value. The height targets are the eight rational abscissa coefficients
of Lemma~\ref{lem:quartic-rational37}.

For every beta-bearing comparison, the common and beta quotient equations
are explicitly
\[
 \Lambda p_H=G_H,\qquad B_A=C_Ap_A,\qquad B_A\Lambda=G_A.
\]
The global nonvanishing in Lemma~\ref{lem:quartic-reciprocal-nonzero}
forces $\Lambda,p_H,B_A,p_A\ne0$. Hence $C_A$ is the ratio of the beta
and common approximate logarithms, with their two minus signs cancelled.
The beta residual is $C_A-2\beta$. In full comparisons $G_H$ has order
$40$ and $G_A$ order $12$; in beta-only comparisons both have order $12$.
The Lucas numerators use the normalized equations below, whose permitted
zero-ordinate branches must not be excluded by this beta argument.

\subsection{Selective comparisons and the local elliptic chart}

\begin{lemma}[Three-mode coverage]\label{lem:quartic-three-mode-coverage}
Number the unchanged raw contacts from $0$ to $36$. Retain full comparisons
at $0,\ldots,4$ and $21,\ldots,28$, beta-only comparisons at
$5,\ldots,12$, and Lucas-only comparisons at $13,\ldots,20$ and
$29,\ldots,36$. The contacts carrying Lucas information cover the three
points $S,T,Z$, and those carrying beta information cover $A,T$.
\end{lemma}
\begin{proof}
The first five positions are shared. Each of $A,S,T,Z$ has its original
block of eight finite positions, in that order. Thus the Lucas set contains
the five shared positions and all three finite blocks of $S,T,Z$; the beta
set contains the five shared positions and the two finite blocks of $A,T$.
The unchanged two-slope rational-height argument applies separately to
these sets. No position is deleted when two slopes, raw values, or assigned
primes coincide. The two maxima used below are consequently
$B_L=\max(|s|,|t|,|z|)$ and $B_{AT}=\max(|n|,|t|)$.
\end{proof}

\begin{lemma}[Doubled formal comparisons]
\label{lem:quartic-doubled-private-comparison}
At a split good prime above the cutoff, if $v(z(R))>0$, then either $R$ or
$R+T_0$ is in the formal group. Whenever the relevant rational function is
defined,
\[
 v(\LogE([2]R))=v(z(R)),
\]
with infinite valuation for $R=O$ or $T_0$.
\end{lemma}
\begin{proof}
For an affine nontorsion point,
$z(R)^2=x(R)/(x(R)^2-D)$. If $v(x)=0$, the right side has nonpositive
valuation. If $v(x)<0$, the point is formal. If $v(x)>0$, translation by
$T_0$ gives $x'=-D/x$ and $z'=-z$, so the translated point is formal with
the same parameter valuation. Doubling kills $T_0$ and is an isometry on
the formal group at an odd good prime. The asserted equality follows from
the retained logarithm isometry. The two other nonzero two-torsion points
are poles of $z$ and are not in the domain.
\end{proof}

Private primary points may therefore lie in $E(K)$ rather than $2E(K)$.
If $[2]P=[2\alpha(P)]P_*$, their coefficients $\alpha(P)$ may be
half-integral. They are integral at each consumed odd prime. One compares
logarithms only after writing an actual integer-multiple identity between
formal points; torsion is not discarded by equating these half-indices.
The primary $f$ tests still put the primary points themselves in the
formal group. An alternate test of
$z(P-Q)/z(J)$, together with the opposite-place $f(J)$ test, implies
$v(z(P-Q))>0$, and the lemma supplies the needed doubled logarithm.

\begin{lemma}[Joint primary and difference graph]
\label{lem:quartic-joint-chart}
Let $Q=(a,aq)$ be nontorsion and $\jmath\ne0$. Set
\[
 F_Q=1/a,\quad c=a+DF_Q,\quad
 r=\jmath/q,\quad s_Q=c/q,\quad t_Q=1/q^2,\quad u=rs_Q.
\]
The three field equations
\begin{align}
 k&=(A-r)p,\\
 kp+up-2Drf-4Dt_QA&=0,\\
 2D(A+r)f+s_Qk+4Dt_QA+\jmath q&=0
\end{align}
give the exact graph $f=1/x(P)$, $p=y(P)/x(P)$,
$A=\jmath z(P-Q)$ for affine primary points $P$ with $x(P)\ne0$
and finite $z(P-Q)$. In particular this primary chart excludes $O$ and
$T_0$. The equations permit $A=0$ and do not
require inverses of $A,k$ or of a difference parameter.
\end{lemma}
\begin{proof}
Set $x=p^2+Df$. The target identity is $c^2-q^4=4D$.
Reversible linear combinations of the displayed equations give
\[
 A U=\jmath H_2,\qquad A V=\jmath H_1,
\]
where
\[
 U=q(x+Df)+cp,\quad V=q^2(x-Df)-4D,\quad
 H_1=q(x+Df)-cp,\quad H_2=x-Df-q^2.
\]
Direct expansion yields
\[
 H_2V-H_1U=-4Dq^2(fp^2+Df^2-1).
\]
Thus $xf=1$ and $P=(x,xp)$ is on $E$. If $U=V=0$, then for
$v\in\{d,-d\}$ one has
$p=-2v/q$ and $x=v(a+v)/(a-v)$, precisely $P=Q+(v,0)$.
At either point $H_2=(4D-q^4)/q^2\ne0$: equality $q^4=4D$ would imply
$N_{K/\Q}(q)^2=28$, impossible for a rational norm. These genuine poles
cannot satisfy the graph equations. Else $U$ or $V$ is nonzero and
$A/\jmath$ is uniquely determined. The chord identities
\[
 H U=N H_2,\qquad H V=N H_1,
\quad H=(p-q)(1+DfF_Q),\quad
 N=pq(1-DfF_Q)+2D(f+F_Q)
\]
identify it with $z(P-Q)$ on the chord domain and hence at every regular
point. At $P=Q$ and $Q+T_0$ the value is zero. Opposition is covered by
$V\ne0$. At $p=0$, the curve has $x=\pm d$, $U=2qx\ne0$, and the graph
still gives its unique regular value. Conversely substituting any point
and its regular difference value reverses these identities and recovers
all three equations with $k=(A-r)p$.
\end{proof}

The last $p=0$ cases are needed below. They are excluded at consumed places
by the positive primary-$f$ valuation, but need not be excluded globally
for every normalized numerator.

\subsection{Private conics, packed field tests and the ordinary pole}

Each contact $r$ has its own rational variables $\theta,s_c,S_c$ with
\[
 S_c=s_c^2,\qquad \theta^2-rS_c=2.
\]
In particular $r,s_c\ne0$. For a field value $F_i=a_i+b_i\sqrt2$ of
required order $e_i$, define
\[
 Y_i=a_i+\theta b_i,\qquad
 T_i=2a_iY_i-a_i^2+2b_i^2
     =Y_i^2-(\theta^2-2)b_i^2,\qquad
 r^{2e_i}W_i=T_i.
\]
The upper tree uses successive pairs $N_r(u,v)=u^2+rv^2$, starting with
the literal leaf $1$ followed by the $W_i$ in stable order, carrying an
unpaired final child unchanged. Denote its root by $W$.

\begin{lemma}[Local conic and tree deductions]
\label{lem:quartic-private-conic-packing}
Assume the ordinary pole gives no negative odd valuation to nonzero
$rW^2$ outside its fixed exceptional primes. At any $p>76176$ with odd
$R=v_p(r)$, the equations above imply $R>0$ and choose one root
$\rho\in\Q_p$ of $2$ such that
$v_p(a_i+\rho b_i)\ge e_i$ for every retained field.
\end{lemma}
\begin{proof}
The conic makes $\delta=v_p(\theta^2-2)=R+2v_p(s_c)$ odd.
It cannot be negative, since then $v_p(\theta^2)$ would be odd; it cannot
be zero. Thus $\delta>0$, $\theta$ is a unit, and one root $\rho$ satisfies
$v_p(\theta-\rho)=\delta$ while the other factor is a unit.
The summands in $N_r(u,v)$ have opposite valuation parity and cannot
cancel. The literal leaf proves $W\ne0$. If $R<0$, ascent along that
leaf gives $v_p(W)\le0$, contradicting the pole conclusion. Hence $R>0$.
The pole conclusion gives $v_p(W)>-R/2$. Descent through the tree uses
\[
 v_p(N_r(u,v))>-cR\ \Longrightarrow
 v_p(u)>-cR/2,\quad v_p(v)>-(c+1)R/2,
\]
with $0\le c<1$, and gives $v_p(W_i)>-R$. Consequently
$v_p(T_i)>(2e_i-1)R\ge2e_i-1$, so $v_p(T_i)\ge2e_i$.
The two summands of $T_i$ also have opposite parity, yielding
$v_p(Y_i)\ge e_i$ and $\delta+2v_p(b_i)\ge2e_i$.
It follows that $v_p((\theta-\rho)b_i)\ge e_i$ and hence the claimed
field order. Zero summands are treated with infinite valuation. No
integrality of $a_i,b_i$ and no global squareclass of $F_i$ was assumed.
\end{proof}

The proof uses $v_p(T_i)\ge2e_i$, not the generally unjustified stronger
bound $2e_iR$. It is independent of the elliptic meaning of the fields.

\begin{lemma}[Conic completeness]\label{lem:quartic-conic-completeness}
In the retained integral-input construction, each contact
$r=d_0^2p_0>0$ admits the needed private conic and arbitrarily deep
approximation to its chosen split branch, independently of the other
contact occurrences.
\end{lemma}
\begin{proof}
Here $p_0\equiv1\pmod8$ and $d_0\ne0$ is a $p_0$-unit. The ring
$\Z[\sqrt2]$ is norm Euclidean: subtracting nearest integers from two
coefficients leaves $|a|,|b|\le1/2$ and $|a^2-2b^2|\le1/2<1$.
The split prime ideal therefore gives integers $a,b$ with $a^2-2b^2=p_0$,
after multiplying by the norm-minus-one unit if needed. Since $b\ne0$,
$(\theta_0,s_0)=(a/b,1/(d_0b))$ is a rational conic point. The parametrization
\[
 \Delta=\frac{2(rs_0-\theta_0 t)}{t^2-r},\qquad
 s_c=s_0+\Delta,\quad \theta=\theta_0+t\Delta
\]
and rational approximation give simultaneous approximation at the finite
set of required places. There is no rational pole $t^2=r$.
Near the selected root, $\delta=1+2v_{p_0}(s_c)$ can be made arbitrarily
large. After all field coefficients are fixed, choose it so that
$\delta+2v_{p_0}(b_i)\ge2e_i$ for every $i$. The original field orders
then make all $W_i$ integral; the constant path makes $W$ a $p_0$-unit.
At the real place the tree is positive, even if some $W_i$ are negative.
The extra approximations used for the following pole may be imposed at
the same time. No conic variable is shared between occurrences.
\end{proof}

\begin{proposition}[Linear-unit ordinary pole]
\label{lem:quartic-linear-pole}
Set $g=3-2\theta$. Replace only the ordinary pole module by
\begin{gather*}
 \ell=gW,\quad z_\ell=\ell^2,\quad c_r=re,\qquad
 p_r^2-17r-34c_r+51ec_r=0,\\
 w=17e-c_rz_\ell,\quad j_s=s_r^2,\quad k_r=rq_r,\\
 j_s^2+34j_sw-867w^2-17q_rk_r=0.
\end{gather*}
It has the one-sided pole conclusion used in
Lemma~\ref{lem:quartic-private-conic-packing}. It has integral-input
completeness under the same inherited realization and rank hypotheses.
The scalar-height pole module is unchanged.
\end{proposition}
\begin{proof}
At an odd prime with odd $v_p(r)$, the conic argument makes $\theta$
integral, and
$(3-2\theta)(3+2\theta)=1-4rs_c^2$ is a unit. Thus $g$ is a unit.
Suppose $b=rW^2\ne0$ has negative odd valuation outside $2,3,17$.
The first conic $p_r^2=17r(1-e)(1+3e)$ forces $e$ to be a unit: positive,
negative or infinite valuation of $e$ would make the right side have odd
valuation. Since $v_p(\ell)=v_p(W)$, one has
$v_p(w)=v_p(b)<0$, an odd integer. In
\[
 (s_r^2-17w)(s_r^2+51w)=17rq_r^2
\]
the left side is nonzero with even valuation, since the valuations of
$s_r^2$ and $w$ differ in parity. The right side is zero or has odd
valuation. This contradiction proves the asserted exclusion, including
zero ordinate branches of the first conic. If $W=0$, the conclusion has
no nonzero pole input; explicit witnesses are
$\ell=z_\ell=p_r=q_r=k_r=0$, $e=1$, $c_r=r$, $w=17$, $s_r=17$, $j_s=289$.

For completeness retain the generalized fixed-$U$ realization argument
under its original hypotheses. For $b>0$ with square reduced denominator
and square in $\Q_2,\Q_3,\Q_{17}$, it applies to a prescribed $U_0$ when
\[
 U_0^2/(17b)<1\text{ at infinity},\quad
 v_2(U_0^2/(17b))\ge3,\quad
 v_3(U_0^2/(17b)),v_{17}(U_0^2/(17b))\ge1,
\]
and $v_p(U_0)=v_p(b)$ at every other odd-support prime of $b$.
Indeed, put $B=17b$ and $V=B(B-U_0^2)$, and let
$\beta_{\rm sf},\eta_{\rm sf}$ be their positive squarefree representatives.
The displayed conditions make $V$ positive and a square in
$\Q_2,\Q_3,\Q_{17}$. Consequently
\[
 \beta_{\rm sf}>1,\quad \beta_{\rm sf}\equiv17\pmod{24},\quad
 \eta_{\rm sf}\equiv1\pmod{24},\quad 17\nmid\eta_{\rm sf}.
\]
At an odd-support prime of $B$ outside $2,3,17$, the square reduced
denominator makes $v_p(b)$ positive, and the hypothesis on $U_0$ gives
$v_p(U_0^2/B)>0$. Thus $1-U_0^2/B$ is a square unit, so
$\eta_{\rm sf}$ is a square in $\Q_p$ and $p\nmid\eta_{\rm sf}$.
The same conclusion at $17$ follows from the displayed smallness condition.
Conversely, at a prime dividing $\eta_{\rm sf}$ the necessary strict
cancellation in $B-U_0^2$ gives $v_p(B)=2v_p(U_0)$ and makes
$B/U_0^2$ a square unit. Thus $\beta_{\rm sf}$ is a square in $\Q_p$.
In particular $\gcd(\beta_{\rm sf},\eta_{\rm sf})=1$, and both
mutual-square conditions hold. If $U_0=0$, then $V=B^2$ and
$\eta_{\rm sf}=1$, so the converse condition is empty. These are precisely
the hypotheses consumed by the inherited simultaneous-realization proof;
no new rank or converse statement is substituted.

Now set $b=rW^2$ and $U_0=rW^2g$. In the complete construction
$r=d_0^2p_0$ and $W$ is a $p_0$-unit, so $v_{p_0}(b)=1$ and all
other valuations of $b$ are even. Hence its reduced denominator is a
square. The retained local squares of $r$ at $2,3,17$ also make $b$ a
square at those places. After the fields are fixed, approximate
the private conic also near $(\theta,s_c)=(3/2,1/(2\sqrt r))$ at infinity,
$2,3,17$. These local points exist by the original contact construction.
The continuous function $(r/17)W(\theta)^2g(\theta)^2$ vanishes there,
so it satisfies the displayed smallness conditions while the selected
prime still makes $W$ a unit. The local-unit argument gives the required
odd-support valuations of $U_0$. The resulting old three-conic witnesses satisfy
\[
 p_o^2=17b(1-e)(1+3e),\quad h_o=e(17b-U_0^2),\quad
 (s_o^2-17bh_o)(s_o^2+51bh_o)=17bq_o^2.
\]
Here $r,W,b$ are nonzero. Set
\[
 p_r=p_o/W,\quad w=h_o/b,\quad
 s_r=s_o/b,\quad q_r=q_o/(r^2W^3).
\]
These substitutions give $w=e(17-r\ell^2)$ and the final conic in the
displayed normalized system; the other coordinates are its definitions.
Thus the witness conversion uses only stated nonzero quantities.
Equivalently the unnormalized form
has $L=r\ell$, $z_L=r^2z_\ell$, $w_L=r^2w$, $s_L=rs_r$,
$q_L=r^2q_r$, $j_L=r^2j_s$, $k_L=r^2k_r$, with $e,p_r,c_r$ unchanged.
All divisions in the reverse scaling are powers of the conic-proved
nonzero $r$. This also proves that normalization loses no zero branch.
\end{proof}

\begin{lemma}[Four-witness anchor]\label{lem:quartic-anchor-four}
For rational $a,b$, both belong to $\Z_{(2)}\cap\Z_{(3)}$ if and only if
$9+3(a^2+b^2)$ is a norm-one trace for the quaternion algebra $(2,3)$.
Consequently the original two-coefficient anchor is equivalent to
\[
 CH-9C^2-3f_0^2-3f_1^2=0,\qquad
 H^2-8U^2-12V^2+24W_a^2-4C^2=0.
\]
\end{lemma}
\begin{proof}
By Definition 4, Observation 5 and Facts (a), (b), (c), (e) in the proof of
Proposition 6 of \cite{koenigsmann-defining}, the algebra is ramified
exactly at $2,3$, the trace set is unrestricted at the other places, and it
is integral at its two ramified places
$2,3$, contains integral traces with residues $1$ at $2$ and $0$ at $3$,
and is determined by these local trace conditions. These are the explicit
trace facts used here, in addition to the older sum-of-traces criterion.
If $m=\min(v_p(a),v_p(b))<0$, then at $3$ the sum of squares has valuation
$2m$, while at $2$ it has valuation $2m$ for unequal valuations and
$2m+1$ for equal finite valuations. Thus $9+3(a^2+b^2)$ would be
nonintegral, proving one implication, with zeros covered by infinite
valuation. Conversely integral $a,b$ give a trace congruent to zero at
$3$. At $2$ it is odd or belongs to $4\Z_2$. For the latter case put
$x_0=h/2\in2\Z_2$, where $h$ is this trace. If $x_0\equiv2\pmod4$, then
$(x_0^2-1)/3\equiv1\pmod8$ is a square, giving
$x_0^2-3v^2=1$. If $x_0\equiv0\pmod4$, then $1-x_0^2/6$ is a square,
and $u=v=1$, $w^2=1-x_0^2/6$ gives
$x_0^2-2u^2-3v^2+6w^2=1$. The local trace criterion supplies rational
witnesses. Finally divide the two homogeneous equations by $C^2$,
which is allowed because the retained constant $C$ is nonzero. This is
an equality of existential projections, not a rational formula choosing
quaternion witnesses.
\end{proof}

\subsection{Local and normalized Lucas tests}

Write
\[
 \mathcal F(s,t,z)=t^2-s(s-2)tz+(s-2)^2z^2-1.
\]
\begin{lemma}[Local Lucas residual]\label{lem:quartic-local-lucas}
The primary and alternate logarithm comparisons, before any proxy equality
is used, give local integral quotients $C_s,C_t,C_z$ congruent to the
integer root indices modulo $p^{40}$. An order-$40$ test of
$\mathcal F(C_s,C_t,C_z)$ implies $p^{40}\mid\mathcal F(s,t,z)$.
\end{lemma}
\begin{proof}
The implication follows by substitution into the polynomial with integer
coefficients. One may evaluate the field residual using
\[
 V=(C_s-2)C_z,\quad D_L=C_sC_t-V,\quad
 E_L=C_t^2-VD_L-1.
\]
These are quadratic equations, and $E_L=\mathcal F(C_s,C_t,C_z)$.
They replace the three proxy comparisons and the global proxy equation.
They do not impose $E_L=0$ as a field equality. Removal of the proxy
coordinates requires their exclusive incidence in that equation and those
comparison arguments. Integer completeness and the reverse input
implication use the two height contradictions below; no conversion of
arbitrary old proxy witnesses is asserted.
\end{proof}

\begin{lemma}[Normalized local Lucas evaluation]
\label{lem:quartic-normalized-lucas}
Let $\Lambda p_H=G_H$ be the common negative-logarithm equation. For the
three Lucas numerators replace their quotient blocks and the preceding
evaluation by
\begin{gather*}
 B_i\Lambda=G_i\quad(i=s,t,z),\qquad P B_s=p_s,\\
 (U+V)p_t=P B_t,\qquad K_Lp_t=(1-2P)B_t,\\
 Vp_z=(1-2P)B_z,\qquad E_L=U^2-K_LV-P^2.
\end{gather*}
An order-$40$ test of $E_L$ has the same required integer consequence.
Every integer input still admits the changed witnesses.
\end{lemma}
\begin{proof}
The global nonvanishing of the retained $G_i$ implies
$\Lambda,p_H,B_s,B_t,B_z\ne0$. Also $p_t\ne0$: otherwise the two
equations with left factor $p_t$ would force both $P=0$ and $1-2P=0$.
At a consumed place all primary points are formal, so their $p_i$ are
nonzero. Put $C_i=B_i/p_i$. Then $P=1/C_s$ and direct substitution gives
\[
 E_L=P^2\mathcal F(C_s,C_t,C_z).
\]
Since $C_s$ is locally integral, the implication
$v(E_L)\ge40\Rightarrow v(\mathcal F(C_s,C_t,C_z))\ge40$ is valid even
when $C_s$ is not a unit. No local denominator has been discarded.

For completeness the integer $s\ne0$ and all root points are fixed before
the prime-pattern theorem is used. Its prescribed finite-exclusion version
allows every chosen prime to avoid the divisors of $s$. Then $C_s$ is a
unit at the consumed place. The old order-$40$ residual implies the new
one, and the displayed equations assign the new witnesses. This imposes
no compatibility between group orders at different primes.

Globally $p_s$ or $p_z$ may vanish at an occurrence not being consumed in
the local argument. Such points have $x=\pm d$ and are covered by the
regular extension in Lemma~\ref{lem:quartic-joint-chart}; they must not be
excluded by reusing an earlier blanket nonzero assertion. At a consumed
place $C_s=2$ gives $P=1/2$ and $V=K_L=0$, since both $p_t$ and $p_z$
are nonzero there. No division by $1-2P$ is needed. Globally the branch
$P=1/2$, $p_z=0$ forces $K_L=0$ but need not force $V=0$; the latter
value may be chosen in the completeness construction.
\end{proof}

The retained field lists can now be stated without proxy coordinates.
In a full occurrence there are the five primary $f$ values, the common
$\overline{f(J)}$, the beta residual, $E_L$, and five alternates.
In a beta-only occurrence there are two primary $f$ values, the common
field, the beta residual and two alternates. In a Lucas-only occurrence
there are four primary $f$ values, the common field, $E_L$ and four
alternates. These lists have lengths $13,6,10$. The formal fields all have
order one. The beta residual and beta alternate have order $12$; the
Lucas residual and Lucas alternates have order $40$. The denominator
alternate has order $e_0=40$ in full and Lucas-only occurrences, and
$e_0=12$ in beta-only occurrences. Its $G$ precision is the same $e_0$.

\begin{theorem}[The two height contradictions]
\label{thm:quartic-selective-arithmetic}
Retain the inherited scalar and contact-height conclusions, and the integer
root indices $n,s,t,z$ of the chosen carrier, with $s,z\ne0$. Let
$\beta=X^2+2$ and $\beta'=2\beta$. Suppose the
local tests above give $C_i\equiv n_i\pmod{p^{e_i}}$ at all consumed
contacts, with beta residual $C_A-\beta'$. Put $H_*=5H$. Suppose
\[
 n_Q^2-6<7H_*N^{19/e},\qquad
 h(\beta')<a/40,\qquad \log(40H_*)<a/40,
 \qquad a=|s|>(24H_*)^{10},
\]
where the first inequality is the contact-height bound whenever the
appropriate radical powers divide $N$. Then the changed arithmetic source
has the same integer-input conclusion as Version 3.
\end{theorem}
\begin{proof}
If $\Delta=\mathcal F(s,t,z)\ne0$, every relevant order-$40$ radical
power divides $|\Delta|$. By three-mode coverage, all three heights in
$B_L=\max(|s|,|t|,|z|)$ satisfy the assumed inequality. Since
$|\Delta|<5B_L^4$, the same elementary estimate as in Version 3 gives
\[
 B_L^{2-76/40}<24H_*.
\]
This contradicts $B_L\ge a>(24H_*)^{10}$. Therefore $\Delta=0$.
The displayed lower bound on $a$ gives $|s|\ge7$, and $z\ne0$ was
retained. The Lucas lemma therefore gives the stronger individual conclusion
$|t|>\exp(a-3)$, not merely a bound on $\max(|t|,|z|)$.

If $\beta'\ne n$, put $N=|\operatorname{num}(\beta'-n)|$.
The beta contacts cover $A,T$, so with $B=\max(|n|,|t|)$ one has
\[
 B^2<7H_*N^{19/12}+6,\qquad
 N\le2e^{h(\beta')}B.
\]
Thus $B^{2-19/12}<40H_*e^{(19/12)h(\beta')}$ and
\[
 \log B<\frac{12+19}{40(24-19)}a=\frac{31}{200}a<a-3,
\]
contrary to $B\ge|t|$. Hence $\beta'=n$ is an integer. $\beta'=2\beta=2(X^2+2)$;
writing $X=u/v$ in lowest terms gives $v^2\mid2$, hence again $v=1$.
\end{proof}


\subsection{General root carrier}
\label{sec:quartic-general-root-variant}

\begin{lemma}[Integer doubled root indices]
\label{lem:quartic-doubled-indices}
Let $R=(1+\sqrt2,2\sqrt2)$, so direct duplication gives $2R=P_*$.
Every $Q\in E(K)$ has a unique integer $n_Q$ with $2Q=n_QP_*$.
It is nonzero exactly when $Q$ is nontorsion. For such a point,
\[
 h(x(Q))>n_Q^2/4-4.
\]
\end{lemma}
\begin{proof}
The inherited generator assertion gives $2Q=n_QP_*$. The torsion
assertion makes its kernel exactly $E[2](K)$, so equivalently
$Q=n_QR+\tau$, $\tau\in E[2](K)$. Uniqueness follows from the infinite
order of $P_*$. The unchanged height bounds give
\[
 n_Q^2-6\le h(x(2Q))\le4h(x(Q))+\log676.
\]
Since $(6+\log676)/4<4$, the stated inequality follows.
\end{proof}

The root points $A,S,T,Z,J_0,J$ may now be any nontorsion points of $E(K)$.
Their indices always mean the integer doubled indices of the lemma. They
are distinct from the half-integer coefficients used for private points
in the earlier carrier convention.

\begin{lemma}[Small general-root charts]
\label{lem:quartic-general-root-charts}
For $A,S,T,Z$ retain $x,y$ and the actual first-doubling square $q=x^2$,
and impose $y^2=xq-Dx$. For $S$ also retain $xf=1$, $yz+x=0$.
Keep the nonzero tangent denominator, or the direct tangent
equation of Lemma~\ref{lem:quartic-tangent-fusion}. For each of $J_0,J$
use $f,z,k$ with
\[
 k=fz,\qquad f=z^2-Dk^2,
\]
with $z\ne0$ enforced by its retained inverse. These charts reconstruct exactly
the required nontorsion root points.
\end{lemma}
\begin{proof}
The first equations are the curve equation with a reused square. The
nonzero tangent denominator excludes all points with $y=0$, hence all
torsion. The two extra S equations preserve its reciprocal and parameter.
In the other chart, $z\ne0$ implies $f\ne0$: $f=0$ would force $k=0$
and then $z=0$. Set $x=1/f$, $p=-1/z$, $y=xp$. Dividing the second
equation by $z^2$ gives $fp^2=1-Df^2$, so $P=(x,y)$ lies on the curve
and $xy\ne0$. Conversely every nontorsion point gives these witnesses.
These are rational maps on the stated nonzero domains. The new equations
cannot justify deleting an unrelated guard without its incidence proof.
\end{proof}

For clarity, the inherited scalar argument still applies to the new
indices. At each odd contact prime of $\beta$ outside $S_*$, the
constant-left scalar norm tree excludes negative odd raw valuation and
gives all six assigned coefficient orders. The first four put $S,J_0$
in the formal group at every place above that prime. Writing
$2S=sP_*$ and $2J_0=j_0P_*$, with nonzero integer $j_0$, the last two
give $v_p(s/j_0)\ge32$, hence $v_p(s)\ge32$. The logarithm comparison
uses the exact identity $[2j_0]S=[2s]J_0$. It is valid also at inert
primes: both rational coefficients are tested and the extension is
unramified outside two. The unchanged scalar raw polynomial and its
small-prime correction detect every prime of $M_5(\beta)$, including
an odd infinite contact. Consequently
\[
 M_5(\beta)^{32}\mid s,\qquad
 h(\beta)\le H M_5(\beta)^{19}\le H|s|^{19/32}.
\]
The height theorem's domain condition holds. The value $\beta=X^2+2$
is none of $0,1,-343/32$, and $\beta=2401$ would require the nonsquare
integer $2399$ to be a rational square. Doubling both indices therefore
changes neither the scalar ratio nor its divisibility argument.

Take $H_*=5H$ and $\beta'=2\beta$. The unchanged contact-height bound
$h(x(Q))<7HN^{19/e}$ and the preceding lemma give
\[
 n_Q^2-6<28HN^{19/e}+7<35HN^{19/e}=7H_*N^{19/e}.
\]
The scalar ratio is still $s/j_0$, since doubling both indices cancels
the factor two. Its order-$32$ tests give
$h(\beta)\le H|s|^{19/32}$. The full anchor at both $2$ and $3$ is
retained. Its $3$-adic consequence, using $2S=sP_*$, is
\[
 s\in2\cdot3^{2^{24}-1}\Z\setminus\{0\},\qquad
 \log_2|s|>25165822.5.
\]
Here formal doubling at $3$ is an isometry and the retained calculation
is $v_3(f(sP_*))=2+2v_3(s)$. The $2$-adic part of the anchor is not
replaced by this consequence. Since $\log_2H_*<2093969$, one has
\[
 |s|>(24H_*)^{10},\qquad |s|>(40H_*)^{32/13},\qquad
 h(2\beta)<|s|/40,\quad\log(40H_*)<|s|/40.
\]
For example the first two required binary cutoffs are respectively
$10\cdot2093974$ and $(32/13)\cdot2093975$, both below the displayed
anchor lower bound. Thus Theorem~\ref{thm:quartic-selective-arithmetic}
applies with unchanged orders. Both beta-bearing comparators must now
test $C_A-2\beta$; the scalar definition $\beta=X^2+2$ and its raw scalar
polynomial remain unchanged.


\subsubsection{Doubled targets}
\label{sec:quartic-target-two-variant}

\begin{lemma}[Targets $2Q_i$ and $P_*$]
\label{lem:quartic-target-two}
The general-root construction may use targets $2Q_i=n_iP_*$ and fixed
denominator target $P_*=2R$, requiring only one doubling of each variable
root. It preserves the local integer congruences and integer completeness.
\end{lemma}
\begin{proof}
For soundness write $2X_0=hP_*$, $2X_i=m_iP_*$ and $2J=j_dP_*$.
The root $J$ is nontorsion, so $j_d\in\Z\setminus\{0\}$.
The doubled differences have indices $h-2$ and $m_i-2n_i$.
Lemma~\ref{lem:quartic-doubled-private-comparison} and the alternate
tests give
\[
 v_p((h-2)/j_d)\ge e_0,\qquad
 v_p((m_i-2n_i)/j_d)\ge e_i,
\]
where $e_0=40$ in full/Lucas mode and $e_0=12$ in beta-only mode.
Since $j_d$ is integral, $h\equiv2$ and $m_i\equiv2n_i$ at those orders.
In particular $h$ is a unit at the odd prime. The exact identity
$[2h]X_i=[2m_i]X_0$ permits logarithm comparison on the primary formal
group, giving ratio $m_i/h$. The retained relative errors, with
$e_i\le e_0$, give $C_i\equiv n_i\pmod{p^{e_i}}$.

For completeness let $N_{p,1},N_{p,2}$ be the two full finite group
cardinalities supplied by the inherited ordinary-prime pattern. It gives
$\gcd(N_{p,1},N_{p,2})=4$ and $p\nmid N_{p,1}N_{p,2}$. Each reduction
has exactly four rational two-torsion points, so its doubling image has
order $M_{p,i}=N_{p,i}/4$. Consequently $M_{p,1}$ and $M_{p,2}$ are
coprime and prime to $p$. Every reduction of a doubled global point is
killed by the corresponding $M_{p,i}$; no global reduction-surjectivity
or cyclicity claim is needed.
Put $j_c=\operatorname{lcm}_p M_{p,2}$ and $J=j_cP_*$, whose doubled
index is $j_d=2j_c$. Separately for each contact choose
\[
 k\equiv0\pmod{M_{p,1}},\quad
 k\equiv1\pmod{M_{p,2}},\quad
 k\equiv1\pmod{p^{40+v_p(j_c)}},\qquad k\notin\{0,1,-1\}.
\]
The three moduli are compatible. Set $X_0=kP_*$ and $X_i=kn_iP_*$.
The primary points are formal at the first place. Their differences from
the chosen targets, and $J$, are formal at the second. The alternate
logarithm ratios are $(k-1)/j_c$ and $(k-1)n_i/j_c$, while the primary
ratios are $n_i$. This supplies every required order. The finite exclusions
avoid zero points and the old open-chart exceptional cases. There is no
simultaneous CRT condition on different contact occurrences.
\end{proof}

The fixed target can be specified directly by
\[
 \begin{array}{c|cc}
 &1&\sqrt2\\\hline
 x(P_*)&57/8&-7/2\\
 q=y(P_*)/x(P_*)&-87/41&129/164\\
 F_Q=1/x(P_*)&456/1681&224/1681\\
 s_Q=(x(P_*)+DF_Q)/q&-55123/14637&19301/58548\\
 t_Q=1/q^2&7624/14161&39904/127449
 \end{array}
\]
Direct substitution verifies the curve and target identities. The source
stores $x,q,F_Q$ and the conjugates of $s_Q,t_Q$. Its remaining coefficients
are affine in $\jmath$:
\[
 r=\jmath\left(-\frac{232}{357}+\frac{86}{357}\sqrt2\right),\qquad
 u=\jmath\left(\frac{291673}{127449}-\frac{29440}{42483}\sqrt2\right),
\]
when $\jmath=1/\overline{z(J)}$. The ten constant bindings and four
affine bindings replace the fixed derivation, subject to full incidence
checks. Their isolated coefficient calculation does not bound the complete
sum of squares.



\subsection{Exact implementation refinements}


\begin{lemma}[Direct tangent equation]\label{lem:quartic-tangent-fusion}
For a curve point with $q=x^2$, replace the inverse and slope equations
$(2y)h=1$, $\lambda=(3q-D)h$ by
$2y\lambda=3q-D$. This preserves all retained rational coordinates.
\end{lemma}
\begin{proof}
The new equation and the curve force $y\ne0$. If $y=0$, then $x=0$ or
$x^2=D$, and $3q-D$ is respectively $-D$ or $2D$, a contradiction.
The unique inverse is restored by $h=(2y)^{-1}$. There is also the
polynomial restoration
\[
 h=\frac{2\lambda(3q-2D)-9xy}{4D^2},
\]
because $(3q-D)(3q-2D)-9x(x^3-Dx)=2D^2$.
For a retained second doubling its input is on the curve: with
$t=X-x$, $Y_0=-y-\lambda t$, the identity
\begin{align*}
 Y^2-X^3+DX={}&y^2-x^3+Dx+t(2y\lambda-3x^2+D)\\
 &+t^2(\lambda^2-X-2x)+(Y-Y_0)(Y+Y_0)
\end{align*}
proves curve preservation from the first doubling equations. Target-two
sources have only the first four variable-root doublings; no second set
of savings or equations is present in that branch.
\end{proof}





\begin{lemma}[Shared field-square evaluation]
\label{lem:quartic-shared-field-square}
For the common field $F=a+b\sqrt2=\overline{f(J)}$, bind once
$c=a^2+2b^2$, $d_s=2ab$. In each designated order-one packed field replace
\[
 Y=a+\theta b,\qquad PW=2aY-a^2+2b^2
\]
by $PW=c+\theta d_s$, with the same $P=r^2$. This is an exact projection
preserving $W$, every higher norm node and the pole witnesses.
\end{lemma}
\begin{proof}
Substitution gives $2a(a+\theta b)-a^2+2b^2=c+\theta d_s$.
The inverse map restores $Y=a+\theta b$, without division. All zero
cases are included. The two coefficients describe $F^2$, but their
packed use still enforces the original order-one condition on $F$;
they must not be relabelled as a new order-one test of $F^2$. The same
identity permits reuse of an already available primary-$f$ square in
its own packed numerator, after checking that square's defining rows.
\end{proof}



\begin{lemma}[First packed denominator and norm fusion]
\label{lem:quartic-first-norm-fusion}
For the first order-one primary field only, replace its virtual denominator
$W_0=T_0/r^2$ by $W_0=T_0/r$, and its first norm $N=1+rW_0^2$ by
\[
 T_0=2a_0Y_0-a_0^2+2b_0^2,\qquad rN-r-T_0^2=0.
\]
Keep every other denominator, the conic, all higher tree nodes and the
ordinary pole. The arithmetic input projection is unchanged.
\end{lemma}
\begin{proof}
The earlier tree proof still gives $v_p(W_0)>-R$, so $v_p(T_0)>0$.
More generally $v_p(T)\ge2e-1$ suffices for an order-$e$ field test:
opposite parity gives $v_p(Y)\ge e$ and
$\delta+2v_p(b)\ge2e-1$, hence
$\delta+v_p(b)\ge e+(\delta-1)/2\ge e$.
Thus the first field still has order one, and every other deduction is
unchanged. For completeness the old depth-two choice has $v_{p_0}(T_0)\ge2$
and $R=1$, so the new virtual leaf is integral and the root remains a unit.
For fixed nonzero $r$, all new tree values remain rational polynomials in
$\theta$, allowing the same exceptional-place pole approximation.

After that arithmetic change, the displayed fusion is exact: set
$T_0=rW_0$ forward and $W_0=T_0/r$ backward. The retained conic proves
$r\ne0$. These maps preserve $N$ and all later nodes, including $N=0$
or $T_0=0$ and negative rational $r$. The unguarded case $r=0$ would be
invalid and is not admitted. The old $r^2$ coordinate remains for other
leaves. The changed field is primary $f(H)$, not the separately shared
common field $\overline{f(J)}$.
\end{proof}


\subsection{Composition, completeness and finite realization}
\label{sec:quartic-composition}

For integral input first set $\beta=X^2+2$ and
take $A=\beta P_*$, whose doubled index is $2\beta$. Choose $J_0,S$ together in the old subgroup with the
required scalar and full anchor divisibilities. Their doubled-index ratio is unchanged. The Lucas existence lemma
then gives $t,z$ for the now fixed $s$, and take $T=tR$,
$Z=zR$. Fix the eight abscissa coefficients and denominator data before
using the prime-pattern theorem. Exclude the divisors of the fixed nonzero
$s$ as well as every earlier finite exceptional set.

Next choose one common comparison point $J$ using the stated least common
multiple, then choose the private points separately for each occurrence,
using precisely the target-two CRT rule.
Construct all joint-graph, reciprocal-$G$, normalized-ratio and
shared-square fields. Only then choose each private conic point and its
norm/pole witnesses by simultaneous approximation. Repeated primes retain
their common orientation but separate conic and private-point witnesses.
The constructions above give every required equation. In the reverse
direction the conic and pole decode field orders at every odd-contact
prime, without an ordinariness assumption. The graph and doubled logarithm
arguments give integer-index congruences, and the two height contradictions
give the integer input. The exact fusions preserve the specified
retained coordinates; arithmetic source replacements use the full
soundness and completeness argument rather than a claim about arbitrary
old private witness tuples.

Finally each selected source must be compiled afresh into rational
quadratic residual equations with the original input coordinate retained.
For real, hence rational, values the all-weights-one sum of their squares
vanishes exactly when every residual vanishes. Clearing nonzero rational
row denominators preserves those zero sets and gives an integral polynomial
of degree at most four.
This elementary last step does not prove the inherited arithmetic inputs.
The exact source, complete flattening map, allowed algebraic transformations,
and full coefficient collection must be checked for the single selected
endpoint. None of the local identities, allocation counts or isolated
coefficient bounds in this appendix is an assertion of that final check.
